% !TEX root = ../thesis.tex

\chapter{Introduction}
\label{ch:introduction}

\section{Background and Motivation}
\label{sec:intro:motivation}

% Migrated and lightly edited from proposal §1 (Introduction) + §2 (Motivation).
The sustained growth of digital health records, together with rapid advances
in artificial intelligence (AI), offers significant potential to transform
healthcare delivery. At the same time, robust regulatory frameworks such as the
General Data Protection Regulation (GDPR) and the Health Insurance Portability
and Accountability Act (HIPAA) play a crucial role in safeguarding patient
confidentiality and ensuring responsible data use. These protections are
essential for maintaining public trust in digital health systems. They also
require researchers to explore innovative methods that enable data-driven
innovation while upholding proper privacy standards. Conventional anonymisation
methods have repeatedly proven insufficient to guarantee privacy in
high-dimensional health data, as re-identification and linkage attacks remain
viable~\cite{Murtaza2023,stadler2022synthetic}.

Synthetic data generation (SDG) has emerged as a promising alternative,
enabling the creation of artificial datasets that retain the statistical
properties of the original while mitigating privacy risks. Among the most
prominent approaches are \emph{Generative Adversarial Networks} (GANs) and
\emph{Large Language Models} (LLMs). GANs have demonstrated success in
structured medical records~\cite{Choi2017MedGAN,Xu2019CTGAN},
and imaging data, while transformer-based LLMs
offer advantages in modelling heterogeneous and unstructured
records~\cite{Miletic2024}. GANs work by training
two neural networks: a generator that creates synthetic samples and a
discriminator that tries to distinguish them from real data until the
generator learns to produce highly realistic outputs. LLMs
operate using transformer architectures that leverage self-attention to learn
long-range patterns in sequences, allowing them to generate coherent text or
structured records. While GANs focus on matching data distributions through
adversarial learning, LLMs generate data by predicting the next most probable
token based on context.

The need for privacy-preserving data generation in healthcare is not only a
technical challenge but a foundational concern of medical research. The
scarcity of openly accessible, high-quality clinical data constrains
innovation in diagnosis, treatment planning, and epidemiological modelling.
Regulatory restrictions and ethical imperatives make direct sharing of
identifiable records infeasible. Synthetic data presents a pathway to alleviate
this bottleneck by allowing the distribution of data without direct disclosure
of personal identifiers, thus accelerating AI model development and
validation.

\section{Problem Statement: The Utility--Privacy Tension}
\label{sec:intro:problem}

The adoption of synthetic data in healthcare depends on two interdependent
requirements: utility and privacy. Synthetic data is valuable only if it remains
analytically useful, and defensible only if it reduces the risk of exposing
information about real patients. The core problem is therefore not to optimise
either criterion in isolation, but to understand the trade-off between them.

\begin{enumerate}\itemsep0.25em
  \item \textbf{Utility:} the dataset must preserve the statistical,
        structural, and predictive characteristics needed for downstream
        analysis, including marginal feature distributions, variable
        dependencies, target class balance, and model generalisation to real
        held-out data.
  \item \textbf{Privacy:} the dataset must provide resistance to adversarial
        attacks, such as membership inference, attribute disclosure, and
        record-level linkage, especially when rare diagnoses, unusual treatment
        patterns, or repeated encounters make patients identifiable even after
        direct identifiers have been removed.
\end{enumerate}
While GANs and LLMs each offer unique strengths, there is limited comparative
work evaluating their effectiveness under a unified privacy--utility framework,
particularly for heterogeneous healthcare datasets.

\section{Research Questions}
\label{sec:intro:rqs}

This thesis addresses three research questions.

\paragraph{RQ1:}
\label{rq:rq1} Which methods and metrics are suitable for evaluating the
utility (statistical resemblance and predictive performance) and privacy risk
of synthetic tabular medical data, and what distinct dimension of data
characteristics does each capture?

\paragraph{RQ2:}
\label{rq:rq2} How do GAN-based and LLM-based synthesisers compare on the
privacy--utility trade-off when applied to tabular medical data, and which
achieves a more effective balance?

\paragraph{RQ3:}
\label{rq:rq3} At a matched privacy budget ($\varepsilon = 5.0$), how does
enforcing differential privacy using mechanisms appropriate to each family
affect utility, and how do the resulting privacy--utility profiles differ
between GAN-based (PATE-GAN vs.\ HealthGAN) and LLM-based (Pythia-1B-DP vs.\
Pythia-1B) synthesisers of tabular medical data?

\section{Contributions}
\label{sec:intro:contributions}

This thesis makes four main contributions to the evaluation of
privacy-preserving synthetic tabular health data.

First, it develops a reproducible evaluation pipeline for synthetic medical
tabular data. The pipeline is organised around two complementary dimensions:
\emph{utility}, comprising both statistical similarity to the real data and the
predictive performance of downstream models, and \emph{privacy}, covering
proximity-based memorisation indicators, membership-inference risk,
attribute-inference risk, and outlier linkability.

Second, it provides a head-to-head empirical comparison of five generators on
the \emph{Diabetes 130-US Hospitals} dataset~\cite{diabetes_db}: HealthGAN
(WGAN-GP)~\cite{healthgan}, PATE-GAN
($(\varepsilon,\delta)$-DP)~\cite{pategan}, Pythia-70M-LoRA and
Pythia-1B-LoRA~\cite{Biderman2023Pythia,Hu2022LoRA,Miletic2024}, and
Pythia-1B-DP-LoRA using DP-SGD via
Opacus~\cite{Yousefpour2021Opacus,Abadi2016DPSGD}.

Third, it quantifies the marginal cost of formal differential privacy within
both model families: non-DP HealthGAN versus PATE-GAN at $\varepsilon=5.0$ in
the GAN setting, and non-DP Pythia-1B versus Pythia-1B-DP at
$\varepsilon=5.0$ in the LLM setting. The inclusion of the non-DP Pythia-1B
scale anchor and the Pythia-1B-DP variant constitutes a scope extension beyond
the originally submitted proposal, motivated by the finding that DP-SGD applied
to Pythia-70M does not converge usefully at this privacy budget
(\S\ref{sec:results:rq3:pythia-70m-failure}).

Fourth, it derives practical recommendations for medical data stewards. These
recommendations identify when a non-DP GAN may be sufficient, which
differentially private mechanism is preferable when a formal guarantee is
required, and under which privacy-budget and governance constraints LLM-based
synthesis is competitive.

\section{Thesis Outline}
\label{sec:intro:outline}

The remainder of this thesis is structured as follows.
\textbf{Chapter~\ref{ch:background}} surveys the background and related work
on tabular synthetic data generation, GAN- and LLM-based approaches,
differential privacy, and existing evaluation frameworks.
\textbf{Chapter~\ref{ch:methodology}} consolidates the experimental
methodology of the thesis. Section~\ref{sec:method:data} describes the
\emph{Diabetes 130-US Hospitals} benchmark and the feature-engineering
pipeline used to prepare it for synthesis.
Section~\ref{sec:method:generators} details the five generators compared in
this work: HealthGAN, PATE-GAN, Pythia-70M-LoRA, Pythia-1B-LoRA, and
Pythia-1B-DP-LoRA, including architectures, training protocols, and the
differentially private mechanisms employed. Section~\ref{sec:method:evaluation} defines the unified
two-pillar evaluation framework: utility (statistical similarity and
predictive performance) and privacy.
\textbf{Chapter~\ref{ch:results}} reports the empirical findings, structured
around the three research questions.
\textbf{Chapter~\ref{ch:discussion}} interprets the results, discusses
strengths and weaknesses of each generator, and acknowledges threats to
validity. Finally, \textbf{Chapter~\ref{ch:conclusion}} summarises the
contributions, returns directly to RQ1--RQ3, and outlines avenues for future
work.

\cleardoublepage
